% !TeX root = ../main.tex
% ---------------------------------------------------------------------------
% Stress, Strain and Constitutive Relations
% Scope extension added from the classical-textbook crosscheck.
% Maps to: Brady & Brown (2006) Ch. 2 "Stress and infinitesimal strain";
% Hoek & Brown (1980) Ch. 5 "Stresses around underground excavations".
% Rationale: the book applies Kirsch, elasto-plastic and finite element
% solutions from Chapter~\ref{ch:ccm} onward without ever establishing strain,
% compatibility or constitutive relations. This chapter supplies that base.
% ---------------------------------------------------------------------------
\chapter{Stress, Strain and Constitutive Relations}
\label{ch:continuum-mechanics}


\section{Why a continuum foundation is needed}
\label{sec:continuum-why}
\scopenote{The Kirsch solution, the convergence--confinement method and every
finite element analysis in this book rest on the same small set of continuum
relations. Stating them once here avoids repeating assumptions later.}

\section{Stress}
\label{sec:stress}

\subsection{Force, traction and the stress tensor}
\label{sec:stress-tensor}
\scopenote{Traction on a plane, the nine components of the stress tensor,
symmetry from moment equilibrium, and the sign convention adopted throughout
this book (compression positive).}

\subsection{Stress transformation and principal stresses}
\label{sec:stress-transformation}
\scopenote{Transformation of axes, the eigenvalue problem giving principal
stresses and directions, and the ordering convention $\sone \ge \stwo \ge
\sthree$.}

\subsection{Stress invariants and the deviatoric split}
\label{sec:invariants}
\scopenote{First, second and third invariants; mean and deviatoric stress;
why yield in rock depends on both the confining and the deviatoric
components.}

\section{Equilibrium and idealisation}
\label{sec:equilibrium-and-idealisation}

\subsection{Equations of equilibrium}
\label{sec:equilibrium}
\scopenote{Differential equilibrium with body forces, in Cartesian and
cylindrical form; the conditions any candidate stress field around an opening
must satisfy.}

\subsection{Plane stress, plane strain and axisymmetry}
\label{sec:plane-idealisations}
\scopenote{The three idealisations that make underground problems tractable,
the physical situations each represents, and the error introduced when a long
tunnel is analysed as plane strain near the face.}

\section{Strain and compatibility}
\label{sec:strain-and-compatibility}

\subsection{Displacement and the strain tensor}
\label{sec:strain-tensor}
\scopenote{Infinitesimal strain from the displacement gradient, engineering
versus tensor shear strain, volumetric and deviatoric strain.}

\subsection{Compatibility}
\label{sec:compatibility}
\scopenote{Why an arbitrary strain field need not correspond to a continuous
displacement field, and the role of compatibility in closed-form solutions.}

\section{Constitutive relations}
\label{sec:constitutive-relations}

\subsection{Elastic constitutive relations}
\label{sec:elastic-constitutive}
\scopenote{Generalised Hooke's law, the relationships between $E$, $\nu$, $G$
and $K$, and the reduced forms for plane stress and plane strain.}

\subsection{Anisotropic elasticity}
\label{sec:anisotropic-elasticity}
\scopenote{Transverse isotropy for bedded and foliated rock, the five
independent constants required, and the practical difficulty of measuring
them.}

\section{Formulations for circular openings}
\label{sec:formulations-for-circular-openings}

\subsection{Cylindrical coordinates for circular openings}
\label{sec:cylindrical-coordinates}
\scopenote{Stress, strain and equilibrium in $r$--$\theta$ form: the natural
setting for the Kirsch solution and for the ground reaction curve of
Chapter~\ref{ch:ccm}.}

\subsection{Elastic strain energy}
\label{sec:strain-energy}
\scopenote{Strain energy density, the energy stored in rock around an
excavation, and its role in the energy-release treatment of
Chapter~\ref{ch:squeezing-rockburst}.}

\section{Inelastic behaviour and graphical methods}
\label{sec:inelastic-behaviour-and-graphical-methods}

\subsection{Beyond elasticity: yield, plasticity and flow}
\label{sec:plasticity-basics}
\scopenote{Yield surface, plastic potential, associated and non-associated
flow, dilation angle, and hardening or softening --- the vocabulary required
for the elasto-plastic solutions and numerical models later in the book.}

\subsection{Graphical representation of stress}
\label{sec:mohr-circle}
\scopenote{The Mohr circle in two and three dimensions, the pole method, and
its use in reading off stress on a discontinuity of given orientation.}
